% !TeX root = closed-systems-from-comparison-completeness.tex
% ============================================================
\chapter{Closed Systems from Comparison Completeness}
\label{ch:closure}
% ============================================================
\begin{chapterorientation}
\orientationitem{Input}{The compactness, comparison-completeness, and quotient-descent results of Chapters~\ref{ch:foundation}--\ref{ch:rotation}.}
\orientationitem{Role}{This chapter consolidates those results into the closure theorem that anchors the rest of the text.}
\orientationitem{Output}{Closure as profile-pairing completeness, together with the minimal semantic package required for representative enrichment.}
\end{chapterorientation}
\section{Introduction}
Under the standing principle of closed-world admissibility
(\cref{sp:closed-world}), this chapter consolidates the
compactness/closure/quotient consequences of
Chapters~\ref{ch:foundation}--\ref{ch:rotation} into a single closure
theorem that serves as the formal entry point to the
representative-enrichment analysis in \cref{ch:locus}.
Chapter~\ref{ch:foundation} isolated the Boolean compactness mechanism
governing passage from finite-stage admissibility to global
admissibility. Chapter~\ref{ch:bottom} applied that mechanism to the
primitive comparison world
\[
	(U,\mathcal C),
\]
proved that rectangular comparison completeness is equivalent to
bijectivity of the canonical factor map
\[
	\Theta:U\to X_A\times X_B,
	\qquad
	X_A:=U/\alpha,
	\qquad
	X_B:=U/\beta,
\]
and thereby obtained the canonical product presentation
\[
	U\cong X_A\times X_B.
\]
Chapter~\ref{ch:rotation} then showed that under this determined product
presentation the intrinsic symmetry group
\[
	G:=\Aut(U,\mathcal C)
\]
acts diagonally on
\[
	X:=X_A\times X_B,
\]
that the comparison world canonically determines the orbit projection
\[
	\pi:X\to\Phys:=X/G,
\]
and that admissible closed-system reports are exactly the coherent maps
on \(X\), equivalently the maps descending uniquely through \(\pi\).
The present chapter does not introduce new primitive structure. Its
purpose is to identify the single internal condition responsible for the
entire Part~II spine and to fix, at theorem level, the meaning of the
phrase \emph{closed system} in the present framework.
From this point onward the convention is exact:
\[
	\text{closed}
	\quad:=\quad
	\text{rectangularly complete}.
\]
This is a terminological identification, not a new axiom. The content of
the chapter is that this one internal completeness condition is exactly
what determines the canonical product presentation, diagonal redundancy, the
orbit quotient, and quotient semantics.
Accordingly, the chapter has three tasks:
\begin{enumerate}[label=\textup{(\roman*)}]
	\item
	      to define closedness exactly as the internal profile-pairing
	      completeness condition isolated in Chapter~\ref{ch:bottom};
	\item
	      to prove that closedness is exactly the condition determining the canonical
	      product presentation and the quotient-semantic structure of
	      Chapter~\ref{ch:rotation};
	\item
	      to prove that closedness is minimal for determining this structure
	      uniformly across comparison worlds.
\end{enumerate}
The first two tasks are consequences of Chapters~\ref{ch:bottom} and
\ref{ch:rotation}. The genuinely new content of the present chapter is
the minimality theorem and the compression of Part~II into a single
classification statement.
Thus this is the point at which the phrase \emph{closed system} ceases
to be heuristic. Within the present framework, a closed system is
exactly a comparison world whose intrinsic left-profile and right-profile
data are internally complete.
% ============================================================
\section{Closure as profile-pairing completeness}
\label{sec:closure-profile}
% ============================================================
Recall from \cref{def:alpha-beta,def:Theta} that every comparison world
\((U,\mathcal C)\) determines intrinsic congruences \(\alpha,\beta\),
quotient sets
\[
	X_A:=U/\alpha,
	\qquad
	X_B:=U/\beta,
\]
and the canonical factor map
\[
	\Theta:U\to X_A\times X_B,
	\qquad
	\Theta(u)=([u]_\alpha,[u]_\beta).
\]
\begin{definition}[Closure]
	\label{def:closure}
	A comparison world \((U,\mathcal C)\) is \emph{closed} if it is
	rectangularly complete in the sense of \cref{def:rect}; equivalently, if
	for every pair
	\[
		A\in X_A,
		\qquad
		B\in X_B,
	\]
	there exists a unique element \(u\in U\) such that
	\[
		[u]_\alpha=A,
		\qquad
		[u]_\beta=B.
	\]
\end{definition}
\begin{remark}[Terminological status]
	\cref{def:closure} introduces no new mathematical structure. It merely
	renames the rectangular comparison completeness condition of
	\cref{def:rect} in the language appropriate to its semantic role in the
	remainder of the monograph.
\end{remark}
\begin{remark}[Internality of closure]
	Closure is purely internal. It refers neither to a geometric boundary
	nor to an ambient environment, an external observer, or a background
	space. It asserts only that every admissible pairing of left and right
	comparison profiles is realized by a unique state of the world itself.
\end{remark}
\begin{proposition}[Closure \(\Longleftrightarrow\) bijectivity of the canonical profile map]
	\label{prop:closure-bijective}
	For a comparison world \((U,\mathcal C)\), the following are equivalent:
	\begin{enumerate}[label=\textup{(\roman*)}]
		\item \((U,\mathcal C)\) is closed.
		\item The canonical factor map
		      \[
			      \Theta:U\to X_A\times X_B
		      \]
		      is bijective.
	\end{enumerate}
	When these conditions hold, the identification
	\[
		U\cong X_A\times X_B
	\]
	is canonical, namely given by \(\Theta\).
\end{proposition}
\begin{proof}
	By \cref{def:closure}, closedness is exactly rectangular completeness.
	By \cref{thm:classification}, rectangular completeness is equivalent to
	bijectivity of \(\Theta\), and the same theorem records that the
	resulting identification of \(U\) with \(X_A\times X_B\) is canonical.
\end{proof}
\begin{remark}[Closure as absence of profile-pair defects]
	Failure of closure means that at least one of two profile-pair
	defects occurs: either \(\Theta\) is not injective, so distinct
	states determine the same left-right profile pair, or \(\Theta\) is
	not surjective, so some left-right profile pair is not realized by
	any state; both defects may occur simultaneously. Thus closure is
	precisely the absence of profile-pair defects.
\end{remark}
% ============================================================
\section{Closure determines diagonal redundancy and quotient semantics}
\label{sec:closure-forces}
% ============================================================
Once closure holds, the product structure and quotient-semantic
architecture of Part~II are determined.
Assume henceforth, under \cref{sp:closure}, that the equivalent conditions of
\cref{prop:closure-bijective} hold. Write
\[
	X:=X_A\times X_B
\]
and identify \(U\) with \(X\) through the canonical bijection
\(\Theta\).
\begin{definition}[Determined diagonal redundancy]
	\label{def:forced-diagonal}
	A comparison world \((U,\mathcal C)\) \emph{determines diagonal redundancy}
	if the following data are canonically determined from \((U,\mathcal C)\)
	without any additional primitive:
	\begin{enumerate}[label=\textup{(\roman*)}]
		\item the factor spaces
		      \[
			      X_A=U/\alpha,
			      \qquad
			      X_B=U/\beta;
		      \]
		\item the canonical identification
		      \[
			      U\cong X_A\times X_B;
		      \]
		\item the canonical diagonal action of
		      \[
			      G:=\Aut(U,\mathcal C)
		      \]
		      on
		      \[
			      X=X_A\times X_B.
		      \]
	\end{enumerate}
\end{definition}
\begin{proposition}[Closure determines diagonal redundancy]
	\label{prop:forced-diagonal-action}
	If \((U,\mathcal C)\) is closed, then it determines diagonal redundancy in
	the sense of \cref{def:forced-diagonal}. More precisely, under the
	canonical identification \(U\cong X_A\times X_B\), the natural action of
	\[
		G=\Aut(U,\mathcal C)
	\]
	on \(U\) becomes the diagonal action
	\[
		g\cdot(x_A,x_B):=(g\cdot x_A,\ g\cdot x_B)
	\]
	on \(X_A\times X_B\).
\end{proposition}
\begin{proof}
	By \cref{prop:closure-bijective}, closure yields the canonical
	identification
	\[
		U\cong X_A\times X_B
	\]
	through \(\Theta\). The \(G\)-invariance of \(\alpha\) and \(\beta\),
	proved in \cref{lem:invariance}, gives well-defined induced actions on
	\(X_A\) and \(X_B\); see
	\cref{def:induced-actions,lem:well-defined-actions}. By
	\cref{thm:diagonal}, the action of \(G\) on \(U\) is transported by
	\(\Theta\) exactly to the diagonal action on \(X_A\times X_B\).
\end{proof}
\begin{corollary}[Canonical quotient input]
	\label{cor:closure-terminal-data}
	If \((U,\mathcal C)\) is closed, then the orbit quotient
	\[
		\Phys:=X/G
	\]
	and the orbit projection
	\[
		\pi:X\to\Phys
	\]
	are canonically determined by the comparison world.
\end{corollary}
\begin{proof}
	By \cref{prop:forced-diagonal-action}, closure determines the diagonal
	action of \(G\) on \(X\). The orbit quotient and its projection are
	therefore canonically determined.
\end{proof}
\begin{proposition}[Universal quotient description of coherent reports]
	\label{prop:principal}
	Under the standing principle \cref{sp:closed-world}, let
	\((U,\mathcal C)\) be closed, and let
	\[
		\pi:X\to \Phys:=X/G
	\]
	be the orbit projection for the determined diagonal action. Then for every
	set \(S\), composition with \(\pi\) induces a bijection
	\[
		\pi^*:\Map(\Phys,S)\xrightarrow{\ \cong\ }\Map_G(X,S),
	\]
	where
	\[
		\Map_G(X,S)
		:=
		\{R:X\to S:\ R(g\cdot x)=R(x)\ \forall g\in G,\ x\in X\}
	\]
	is the set of coherent, equivalently \(G\)-invariant, maps.
\end{proposition}
\begin{proof}
	Define
	\[
		\pi^*(f):=f\circ\pi
		\qquad
		\text{for }f\in\Map(\Phys,S).
	\]
	Because \(\pi(g\cdot x)=\pi(x)\) for all \(g\in G\) and \(x\in X\), the
	composite \(f\circ\pi\) is \(G\)-invariant. Hence \(\pi^*\) is well
	defined.
	Let
	\[
		R\in\Map_G(X,S).
	\]
	Then \(R\) is coherent in the sense of \cref{def:report-coherence}. By
	\cref{thm:descent}, there exists a unique map
	\[
		\widetilde R:\Phys\to S
	\]
	such that
	\[
		R=\widetilde R\circ\pi.
	\]
	Therefore
	\[
		R=\pi^*(\widetilde R),
	\]
	so \(\pi^*\) is surjective.
	Now suppose
	\[
		\pi^*(f_1)=\pi^*(f_2),
		\qquad\text{i.e.}\qquad
		f_1\circ\pi=f_2\circ\pi.
	\]
	Let \(p\in\Phys\). Since \(\pi\) is surjective, choose \(x\in X\) with
	\[
		\pi(x)=p.
	\]
	Then
	\[
		f_1(p)
		=
		f_1(\pi(x))
		=
		(f_1\circ\pi)(x)
		=
		(f_2\circ\pi)(x)
		=
		f_2(\pi(x))
		=
		f_2(p).
	\]
	Hence \(f_1=f_2\), so \(\pi^*\) is injective. Therefore \(\pi^*\) is a
	bijection.
\end{proof}
\begin{corollary}[Closed-system semantics is quotient semantics]
	\label{cor:closed-equals-quotient}
	Under closure, admissible state reports in the closed-system sense are
	exactly maps on the quotient state space
	\[
		\Phys=X/G.
	\]
\end{corollary}
\begin{proof}
	By \cref{def:closed-system-semantics}, admissible state reports are
	exactly coherent maps, equivalently the elements of \(\Map_G(X,S)\) for
	varying codomain \(S\). By \cref{prop:principal}, these are in canonical
	bijection with maps \(\Phys\to S\).
\end{proof}
\begin{remark}[Why quotient semantics is not an extra axiom]
	The quotient \(\Phys\) is not introduced by a separate gauge principle.
	Once closure determines the canonical product structure and diagonal
	redundancy, the universal property of \(\pi\) determines the semantic
	identification of admissible reports with maps on \(\Phys\).
\end{remark}
% ============================================================
\section{Logical minimality of closure}
\label{sec:closure-minimality}
% ============================================================
The genuinely new content of the chapter is the minimality statement:
closure is not merely sufficient for the quotient-semantic spine. It is
the minimal internal condition that determines it uniformly.
\begin{theorem}[Logical minimality of closure]
	\label{thm:logical-minimality}
	Let \(P\) be any property of comparison worlds with the following
	uniform determining property:
	\medskip
	\noindent
	for every comparison world \((U,\mathcal C)\) satisfying \(P\), the
	world canonically admits an identification
	\[
		U\cong X_A\times X_B
	\]
	compatible with the intrinsic congruences \(\alpha\) and \(\beta\).
	\medskip
	\noindent
	Then \(P\) implies rectangular completeness. Equivalently, \(P\) implies
	closedness.
\end{theorem}
\begin{proof}
	Fix a comparison world \((U,\mathcal C)\) satisfying \(P\). By
	hypothesis, \((U,\mathcal C)\) canonically admits an identification
	\[
		U\cong X_A\times X_B
	\]
	compatible with the intrinsic congruences \(\alpha\) and \(\beta\). By
	compatibility, the canonical factor map
	\[
		\Theta:U\to X_A\times X_B
	\]
	must coincide with that identification, hence is bijective. By
	\cref{thm:classification}, bijectivity of \(\Theta\) is equivalent to
	rectangular completeness. Thus \(P\) implies closedness.
\end{proof}
\begin{theorem}[Minimality of closure for the quotient-semantic spine]
	\label{thm:rect-minimal}
	For a comparison world \((U,\mathcal C)\), the following are equivalent:
	\begin{enumerate}[label=\textup{(\arabic*)}]
		\item \((U,\mathcal C)\) is closed.
		\item The canonical map
		      \[
			      \Theta:U\to X_A\times X_B
		      \]
		      is bijective.
		\item \((U,\mathcal C)\) determines diagonal redundancy.
	\end{enumerate}
	Moreover, no strictly weaker closure property guarantees determined diagonal
	redundancy for all comparison worlds.
\end{theorem}
\begin{proof}
	The equivalence
	\textup{(1)}\(\Longleftrightarrow\)\textup{(2)} is
	\cref{prop:closure-bijective}. The implication
	\textup{(2)}\(\Longrightarrow\)\textup{(3)} is
	\cref{prop:forced-diagonal-action}. Conversely, if diagonal redundancy is
	determined, then by \cref{def:forced-diagonal} the comparison world
	canonically determines an identification
	\[
		U\cong X_A\times X_B.
	\]
	Since \(\Theta\) is the canonical map from \(U\) to that same product,
	it must be bijective. Hence
	\textup{(3)}\(\Longrightarrow\)\textup{(2)}.
	If \(P\) is any property of comparison worlds which uniformly determines
	diagonal redundancy, then \(P\) satisfies the hypothesis of
	\cref{thm:logical-minimality}. Therefore \(P\) implies closedness. Hence
	no strictly weaker property can force diagonal redundancy uniformly.
\end{proof}
\begin{corollary}[Logical irreversibility of the spine]
	Once closure holds, the structural chain
	\begin{align*}
		(U,\mathcal C)
		 & \Longrightarrow U\cong X_A\times X_B                     \\
		 & \Longrightarrow G\curvearrowright X \text{ diagonally}   \\
		 & \Longrightarrow \pi:X\to\Phys=X/G                        \\
		 & \Longrightarrow \text{canonical descent}                    \\
		 & \Longrightarrow \text{subsystem attribution obstruction}
	\end{align*}
	is determined. No weaker closure axiom guarantees this chain uniformly.
\end{corollary}
\begin{proof}
	The first three steps follow from
	\cref{prop:closure-bijective,prop:forced-diagonal-action,cor:closure-terminal-data}.
	Canonical descent is the universal quotient description of
	\cref{prop:principal}. The final obstruction statement is
	\cref{thm:nogo}. The minimality clause is the second assertion of
	\cref{thm:rect-minimal}.
\end{proof}
% ============================================================
\section{Part~II classification theorem}
\label{sec:partII-classification}
% ============================================================
The results of Chapters~\ref{ch:bottom}, \ref{ch:rotation}, and the
present chapter now compress into a single classification theorem.
\begin{theorem}[Closed-system classification from comparison completeness]
	\label{thm:structural-closure}
	Under the standing principle \cref{sp:closed-world}, let
	\((U,\mathcal C)\) be a comparison world in the local
	distinguishability regime
	\[
		\alpha\cap\beta=\Delta_U
	\]
	and satisfying the closure clause \cref{sp:closure}, equivalently
	rectangular comparison completeness.
	Then the following structure is canonically determined.
	\begin{enumerate}[label=\textup{(\roman*)}]
		\item
		      the canonical factor map
		      \[
			      \Theta:U\to X_A\times X_B,
			      \qquad
			      X_A:=U/\alpha,
			      \qquad
			      X_B:=U/\beta,
		      \]
		      is bijective, hence
		      \[
			      U\cong X_A\times X_B;
		      \]
		\item
		      the intrinsic symmetry group
		      \[
			      G:=\Aut(U,\mathcal C)
		      \]
		      acts diagonally on
		      \[
			      X:=X_A\times X_B,
		      \]
		      so the comparison world canonically determines the orbit projection
		      \[
			      \pi:X\to\Phys:=X/G;
		      \]
		\item
		      for every set \(S\), pullback along \(\pi\) induces a bijection
		      \[
			      \Map(\Phys,S)\xrightarrow{\ \cong\ }\Map_G(X,S),
		      \]
		      where \(\Map_G(X,S)\) denotes the coherent, equivalently \(G\)-invariant,
		      maps \(X\to S\);
		\item
		      consequently, admissible closed-system content is exactly quotient
		      content on \(\Phys\);
		\item
		      for pure orbit motion, coherent reports are constant, so subsystem
		      attribution along pure diagonal orbit loops is not a quotient-level
		      invariant;
		\item
		      any extension of quotient semantics on finite comparison protocols that
		      is not endpoint-determined introduces additional structure through at
		      least one of two loci, possibly both, and through no others:
		      \begin{enumerate}[label=\textup{(\alph*)}]
			      \item object-level representative choice, i.e. a section
			            \[
				            s:\Phys\to X;
			            \]
			      \item morphism-level transport data, equivalently nontrivial loop defect.
		      \end{enumerate}
		\end{enumerate}
		Hence a closed system in the present sense is canonically classified by
		the quotient-semantic datum
		\[
			X_A,\ X_B,\ G,\ \pi:X_A\times X_B\to (X_A\times X_B)/G,
		\]
		together with the theorem that all admissible content descends to the
		quotient and all non-quotient enrichment enters through at least one of
		the two classified loci, possibly both, with no third locus.
\end{theorem}
\begin{proof}
	Item \textup{(i)} is \cref{prop:closure-bijective}. Item \textup{(ii)}
	is \cref{prop:forced-diagonal-action} together with
	\cref{cor:closure-terminal-data}. Item \textup{(iii)} is
	\cref{prop:principal}. Item \textup{(iv)} is
	\cref{cor:closed-equals-quotient}. Item \textup{(v)} is
	\cref{thm:nogo}. Item \textup{(vi)} is
	\cref{thm:two-locus-classification,cor:transport-form}.
\end{proof}
% ============================================================
\section{Interpretation}
\label{sec:closure-interpretation}
% ============================================================
The formal results admit an exact internal reading.
\begin{remark}[Closure as internal completeness]
	Within the present framework, the term \emph{closed system} has a purely
	comparison-theoretic meaning. It does not refer to a geometric boundary
	or to an externally imposed isolation condition. Rather, closure means
	that the comparison structure already realizes every state required by
	its own left-profile and right-profile data.
	Equivalently, closure is the condition that every pair consisting of an
	\(\alpha\)-class and a \(\beta\)-class is realized by a unique element
	of \(U\). This is exactly the rectangular comparison completeness
	condition of \cref{def:rect}, recast terminologically in
	\cref{def:closure}.
	Thus closedness is not added from outside. It is the statement that the
	comparison world is internally complete with respect to its own profile
	data. Once this internal completeness holds, the canonical product
	decomposition, the diagonal symmetry, the orbit quotient, and the
	quotient semantics all follow as necessary consequences of the
	comparison structure itself.
\end{remark}
\begin{remark}[Structural centrality of the chapter]
	Everything later in the stack presupposes the theorem-level content
	recorded here. The transport theory of enrichment, the stabilization of
	loop defect, and the later curvature realization of that same
	obstruction all require that admissible state content has already
	been determined down to quotient semantics. That determining is exactly
	the content of closure as established in the present chapter.
\end{remark}
% ============================================================
\section{Conclusion}
\label{sec:closure-conclusion}
% ============================================================
The structural content of Part~II is now complete.
A comparison world is closed exactly when it satisfies rectangular
comparison completeness. Equivalently, it is closed exactly when the
canonical factor map
\[
	\Theta:U\to X_A\times X_B
\]
is bijective. In that case the comparison structure determines the canonical
product presentation
\[
	U\cong X_A\times X_B,
\]
the diagonal action of
\[
	G=\Aut(U,\mathcal C)
\]
on
\[
	X=X_A\times X_B,
\]
and hence the canonical orbit projection
\[
	\pi:X\to\Phys=X/G.
\]
The universal property of \(\pi\) then identifies admissible
closed-system reports with maps on \(\Phys\). Quotient semantics is
therefore not an added axiom. It is a theorem-level consequence of
closure.
	Moreover, closure is minimal for this purpose. No strictly weaker
	condition determines diagonal redundancy uniformly, and hence no strictly
	weaker condition determines the quotient-semantic spine of Part~II.
	This is the structural role of the chapter. It closes Part~II by fixing
	the exact theorem-level meaning of a closed system and by proving that
	the passage from primitive comparison data to canonical quotient
	semantics is determined exactly by closure, i.e. by rectangular
	comparison completeness.
	Accordingly, Part~II concludes at theorem level: closure is both
	necessary and sufficient for the quotient-semantic backbone, and no
	weaker hypothesis supports the same uniform determination.

From this fixed backbone, \cref{ch:locus} undertakes the universal
classification of admissible enrichment mechanisms beyond quotient
semantics.

\begin{chapterclosing}
\closingitem{Established}{Closure is exactly rectangular comparison completeness, and it is necessary and sufficient for the quotient-semantic backbone.}
\closingitem{Carried forward}{The remaining question is not the quotient itself, but the classification of admissible enrichment beyond it.}
\end{chapterclosing}
